% BHU PYQ -- Philosophy: Western Philosophy (2024-25 NEP)
\documentclass[11pt,a4paper]{article}
\usepackage[a4paper,top=2.5cm,bottom=2.5cm,left=2.8cm,right=2.8cm,headheight=14pt]{geometry}
\usepackage{amsmath,amssymb}
\usepackage{fontspec}
\setmainfont{Kohinoor Devanagari}
\usepackage{microtype,fancyhdr,enumitem,hyperref,lastpage,parskip}

\hypersetup{colorlinks=true,urlcolor=black,linkcolor=black,pdftitle={PHIMJ21: Western Philosophy -- BHU NEP 2024-25}}

\pagestyle{fancy}\fancyhf{}
\renewcommand{\headrulewidth}{0.4pt}\renewcommand{\footrulewidth}{0.4pt}
\fancyhead[L]{\small \textit{Banaras Hindu University}}
\fancyhead[R]{\small \textit{PHIMJ21: Western Philosophy (2024-25)}}
\fancyfoot[C]{\small Page \thepage\ of \pageref{LastPage}}

\newlist{parts}{enumerate}{2}
\setlist[parts,1]{label=(\alph*),leftmargin=*,itemsep=4pt,topsep=3pt}
\setlist[parts,2]{label=(\roman*),leftmargin=*,itemsep=2pt,topsep=2pt}

\newcommand{\pts}[1]{\hfill \textit{[#1]}}

\begin{document}

\begin{center}
  {\large\bfseries BANARAS HINDU UNIVERSITY}\\[2pt]
  {\normalsize B. A. (Hons.) Semester II Examination, 2024-25 (NEP)}\\[6pt]
  \rule{0.8\linewidth}{0.5pt}\\[6pt]
  {\large\bfseries PHILOSOPHY}\\[4pt]
  {\large\bfseries Paper Code: PHIMJ-21 : Western Philosophy}\\[4pt]
  {\normalsize Time: 2:30 Hours \hfill Full Marks: 70}\\[2pt]
  {\small REF NO. CE/Conf./D-II/38}
\end{center}
\rule{\linewidth}{0.4pt}
\smallskip

\noindent \textit{Instructions: Attempt all Sections A, B and C according to instructions given. Write your Roll No.\ at the top immediately on receipt of this question paper.}

\smallskip
\noindent \textit{दिये गये निर्देशों के अनुसार खण्ड अ, ब तथा स के उत्तर दीजिये।}

\bigskip

\section*{SECTION -- A / खण्ड -- अ}
\noindent \textit{Note: Attempt any \textbf{three} questions, each in 300 words. Each question carries 10 marks.}\pts{3 $\times$ 10 = 30}

\noindent \textit{किन्हीं तीन प्रश्नों के उत्तर दीजिये, प्रत्येक उत्तर 300 शब्दों में हो। प्रत्येक प्रश्न 10 अंकों का है।}

\begin{enumerate}[label=\textbf{\arabic*.},leftmargin=*,itemsep=12pt]

  \item Explain Descartes' Mind-Body dualism and discuss how he reconciles them.\pts{10}
  
  डेकार्ट के मन-शरीर द्वैतवाद की व्याख्या कीजिये तथा उसकी विवेचना कीजिये।

  \item Discuss Spinoza's doctrine of Pantheism in detail.\pts{10}
  
  स्पिनोजा के सर्वेश्वरवाद के सिद्धान्त की विस्तार से चर्चा कीजिये।

  \item Critically examine Locke's rejection of Innate Ideas.\pts{10}
  
  जन्मजात प्रत्ययों के लॉक द्वारा खण्डन का समालोचनात्मक परीक्षण कीजिये।

  \item What is Hume's Skepticism? How does he refute the principle of causality?\pts{10}
  
  ह्यूम का सन्देहवाद क्या है? वे कारणता के सिद्धान्त का खण्डन किस प्रकार करते हैं?

  \item How does Kant reconcile Rationalism and Empiricism in his Critical Philosophy?\pts{10}
  
  काण्ट अपने समीक्षात्मक दर्शन में बुद्धिवाद तथा अनुभववाद का समन्वय किस प्रकार करते हैं?

\end{enumerate}

\bigskip

\section*{SECTION -- B / खण्ड -- ब}
\noindent \textit{Note: Attempt any \textbf{four} questions out of the following six, each in 150 words. Each question carries 5 marks.}\pts{4 $\times$ 5 = 20}

\noindent \textit{किन्हीं चार प्रश्नों के उत्तर दीजिये, प्रत्येक उत्तर 150 शब्दों में हो। प्रत्येक प्रश्न 5 अंकों का है।}

\begin{enumerate}[label=\textbf{\arabic*.},leftmargin=*,start=6,itemsep=10pt]

  \item Throw light on the concept of Substance according to Spinoza.\pts{5}
  
  स्पिनोजा के अनुसार द्रव्य (Substance) की अवधारणा पर प्रकाश डालिये।

  \item Discuss Berkeley's Subjective Idealism (Esse est percipi).\pts{5}
  
  बर्कले के आत्मनिष्ठ प्रत्ययवाद की चर्चा कीजिये।

  \item What is Hegel's notion of Dialectic?\pts{5}
  
  हीगेल का द्वन्द्वात्मक सिद्धान्त क्या है?

  \item When Descartes says `Cogito, ergo sum' (I think, therefore I am), what does he mean?\pts{5}
  
  जब डेकार्ट कहते हैं 'मैं सोचता हूँ, अतः मैं हूँ', तो उनका क्या अर्थ है?

  \item Explain what Kant meant by calling his approach ``Critical Philosophy''.\pts{5}
  
  काण्ट द्वारा अपनी पद्धति को "समीक्षात्मक दर्शन" कहने के महत्त्व की व्याख्या कीजिये।

  \item Explain Leibniz's theory of Monadology in short.\pts{5}
  
  लाइब्निट्ज़ के चिद्णुवाद (Monadology) के सिद्धान्त को संक्षेप में समझाइये।

\end{enumerate}

\bigskip

\section*{SECTION -- C / खण्ड -- स}
\noindent \textit{Note: Attempt all \textbf{ten} questions of this section in a maximum of two sentences each. Each question carries 2 marks.}\pts{10 $\times$ 2 = 20}

\noindent \textit{इस खण्ड के सभी दस प्रश्नों के उत्तर अधिकतम दो वाक्यों में दीजिये। प्रत्येक प्रश्न 2 अंकों का है।}

\begin{enumerate}[label=\textbf{\arabic*.},leftmargin=*,start=12,itemsep=8pt]
  \item 
  \begin{parts}
    \item Who is called the Father of Modern Western Philosophy?\pts{2}
    
    आधुनिक पाश्चात्य दर्शन का जनक किसे कहा जाता है? (René Descartes)

    \item Name the book written by Baruch Spinoza on Ethics.\pts{2}
    
    स्पिनोजा द्वारा नीतिशास्त्र पर लिखी गयी पुस्तक का नाम लिखिये। (Ethics)

    \item What is the main thesis of John Locke's Empiricism?\pts{2}
    
    जॉन लॉक के अनुभववाद का मुख्य सिद्धान्त क्या है? (Tabula Rasa)

    \item What does `Esse est percipi' mean?\pts{2}
    
    `सत्ता अनुभवमूलक है' (Esse est percipi) का क्या अर्थ है?

    \item Name the two kinds of perceptions according to David Hume.\pts{2}
    
    डेविड ह्यूम के अनुसार प्रत्यक्ष के दो प्रकारों के नाम लिखिये। (Impressions and Ideas)

    \item Which philosopher wrote ``Critique of Pure Reason''?\pts{2}
    
    ``क्रिटिक ऑफ प्योर रीजन'' (शुद्ध बुद्धि की समीक्षा) ग्रन्थ किसने लिखा? (Immanuel Kant)

    \item What is a Monad according to Leibniz?\pts{2}
    
    लाइब्निट्ज़ के अनुसार चिद्णु (Monad) क्या है?

    \item What is Pre-established Harmony?\pts{2}
    
    पूर्व-स्थापित सामंजस्य का सिद्धान्त क्या है?

    \item Define Synthetic A Priori Judgments according to Kant.\pts{2}
    
    काण्ट के अनुसार संश्लेषणात्मक प्रागुभवात्मक निर्णय को परिभाषित कीजिये।

    \item Who propounded the Absolute Idealism in 19th Century German philosophy?\pts{2}
    
    19वीं शताब्दी के जर्मन दर्शन में निरपेक्ष प्रत्ययवाद का प्रतिपादन किसने किया? (G. W. F. Hegel)
  \end{parts}
\end{enumerate}

\end{document}
